\documentclass{article}

\usepackage{microtype}
\usepackage{graphicx}
\usepackage{subcaption}
\usepackage{booktabs}
\usepackage{tabularx}
\usepackage{hyperref}
\usepackage{amsmath}
\usepackage{amssymb}
\usepackage{multirow}
\usepackage{array}
\usepackage{xcolor}
\usepackage{balance}

\newcommand{\theHalgorithm}{\arabic{algorithm}}
\newcommand{\norm}[1]{\left\lVert #1 \right\rVert}
\newcommand{\R}{\mathbb{R}}
\newcommand{\od}{\mathbin{\odot}}
\newcommand{\DAWN}{DAWN-SRW}
\newcommand{\smallhead}[1]{\vspace{0.25em}\noindent\textbf{#1}}
\newcolumntype{Y}{>{\raggedright\arraybackslash}X}

\usepackage[preprint]{icml2026}
\makeatletter
\renewcommand{\Notice@String}{}
\makeatother
\hypersetup{pdfsubject={}}

\icmltitlerunning{\smash{DAWN-SRW: State-Conditioned Programs}}

\begin{document}

\twocolumn[
\icmltitle{DAWN-SRW: State-Conditioned Programs for Learned State Transitions}

\begin{icmlauthorlist}
\icmlauthor{Seungho Choi}{ind}
\end{icmlauthorlist}

\icmlaffiliation{ind}{Independent Researcher}
\icmlcorrespondingauthor{Seungho Choi}{madst0614@gmail.com}
\icmlkeywords{State-Conditioned Programs, Language Modeling, Learned State Transitions, Read-Write Operators, Conditional Computation, Mechanistic Analysis}

\vskip 0.3in
]

\printAffiliationsAndNotice{}

\begin{abstract}
Learned state-transition computation is often organizationally implicit. We ask
whether it can be organized so that the conditions, operations, and transition
effects constituting its execution are explicit objects of measurement and
intervention. We study this question constructively in autoregressive language
models. We use \((\mathcal O,C,\Gamma)\) as minimal organizational coordinates
for one family of explicit learned transition organizations. Within these coordinates,
a \emph{State-Conditioned Program} is one program-oriented formulation in which
controlled operation contributions are constituents of actual forward
execution. It is not a claim of symbolic program text, formal operational
semantics, or an intrinsic or privileged neural decomposition. \DAWN{} provides one constructive
realization using rank-1 read-write (RW) operators that read a scalar from the
current representation and write along a learned transition direction, together
with an operation-address space, normalized additive composition, full candidate
scope, and dense RW execution. At 393.8M parameters and 39.999B tokens, \DAWN{} supports
viability relative to the matched dense control: C4 cross-entropy is 2.9402
versus 2.9359 and the six-task zero-shot mean is 0.4629 versus 0.4597.
State-conditioned exact applicability support provides evidence of conditional
organization. A discovery-frozen 4,096-site Q/K-centered operator circuit passes
suppression, four matched-control families, and exact restoration on validation
and held-out indirect-object-identification (IOI) inputs, establishing causal
accessibility; fail-closed results delimit unsuccessful localization. They do
not establish human-readable operator semantics, benchmark superiority,
hardware efficiency, or an intrinsic or unique decomposition.
\end{abstract}

\section{Introduction}
\label{sec:introduction}

Learned systems compute through state transitions, but their transition-producing
organization is often implicit. State dependence and conditional structure are
established; the broader problem is whether learned computation can be
architecturally organized around explicit conditions, operations, and transition
effects. Autoregressive language models provide a concrete test setting in which
high-dimensional representations undergo repeated transitions. If
\(X_\ell\) denotes the sequence
representation at layer \(\ell\), a block implements, at a coarse level,
\(X_\ell\mapsto X_{\ell+1}\).

We therefore ask:
\begin{quote}
\emph{Can learned state-transition computation be organized so that the
conditions, operations, and transition effects constituting its execution are
explicit objects of measurement and intervention?}
\end{quote}

Autoregressive language modeling is the present constructive test domain, not a
boundary on this research question.

At the broadest level, program execution and dynamical systems share the
substrate
\(x\longrightarrow x'\): the former treats computation as state or
configuration change, while the
latter treats the same abstract structure as state evolution. This commonality
does not give \DAWN{} a formal operational semantics. A more specific
organizational motif,
\(
\text{condition}\longrightarrow\text{operation}\longrightarrow
\text{state transition},
\)
is likewise old; we do not claim it as novel. To separate its descriptive
roles, we use the organizational coordinates
\begin{equation}
\mathcal P=(\mathcal O,C,\Gamma),
\label{eq:organizational_coordinates}
\end{equation}
where \(\mathcal O=\{O_i\}_i\) is an operation family, \(C\) is the
state-conditioned applicability relation, and \(\Gamma\) is the composition
rule. For scalar applicability,
\begin{equation}
C(x,O_i)=c_i(x),\qquad
F_{\mathcal P}(x)=\Gamma\!\left(\{c_i(x)O_i(x)\}_i\right).
\label{eq:general_composition}
\end{equation}
These minimal coordinates describe one family of explicit learned transition
organizations, not a universal ontology or necessary decomposition of learned
dynamics. The form does not itself state the architectural hypothesis.
Program and execution traditions motivate identifying conditioned operations;
dynamics and control provide the richer mathematical precedent for continuously
valued validity and composition of overlapping component effects. The
coordinates distinguish what transformations exist, when and how much they
apply, and how their controlled contributions form a transition-producing
result. Within them, the \emph{control--operation factorization} separates
\(C\) from \(\mathcal O\); \(\Gamma\) is a separate composition role, not a
third factor of that factorization. Normalized additive composition is one
\DAWN{} realization of \(\Gamma\), not part of the coordinates' definition.

The coordinates alone are not the scientific claim. An arbitrary mapping
\(T(x)\) is vacuously compatible by taking one operation \(O_1=T\) and
\(c_1(x)=1\). Such post-hoc functional decomposability is weaker than a learned
system whose forward execution is actually organized through operation
identities, state-conditioned applicability, and controlled contributions; a
further architectural step exposes those constituents as explicit variables or
objects. Our interest is in this computational organization and architectural
exposure, not in algebraic decomposability by itself.

Within this broader problem, we call one program-oriented organizational
formulation a \emph{State-Conditioned Program}: the current state determines
operation applicability, and controlled operation contributions are constituents
of the actual transition-producing computation. This definition requires neither
continuous or geometric control, sparse support, rank-1 operations, persistence,
nor additive composition. The term \emph{program} generalizes the analogy beyond
symbolic source, discrete instructions, named human-readable primitives,
call/return control flow, or formal operational semantics.

The same learned computation admits complementary descriptions. A neural-network
view asks what tensor map \(x\mapsto F_{\mathcal P}(x)\mapsto x'\) is implemented;
a dynamics view asks what evolution \(x\mapsto x'\), or residual update
\(x\mapsto x+\Delta x\), occurs. The State-Conditioned Program view asks which
operations contribute at a state, under what conditions, by how much, and what
happens when selected contributions are changed. It is therefore a computational
view grounded in actual execution, not an ontological claim that this is the only
true, intrinsic, necessary, or privileged neural decomposition. Operational
accounts of transitions between configurations provide an operation-oriented
conceptual lineage \cite{plotkin1981structural}, but no formal semantics for
\DAWN{} is claimed.

The resulting hierarchy is conceptual, not chronological. Its convergence point
is explicitly organized learned state-transition computation, not the
State-Conditioned Program formulation itself. Dynamics and control provide strong
mathematical precedent for graded state-dependent validity, local transition
models, and weighted composition
\cite{zadeh1965fuzzy,takagi1985fuzzy,johansen1995identification,
rugh2000research}. Learned conditional and modular systems then make both
\(C_\theta\) and \(O_{\theta,i}\) learned computational objects
\cite{shazeer2017outrageously,rosenbaum2018routing,goyal2021rim,
rahaman2021dynamic}; stable operation identity is an architectural property,
whereas reuse is an empirical possibility. Read--write and fast-weight work
supplies an elementary transformation structure \cite{schlag2021linear,
lu2026little}. The architectural hypothesis studied here is that learned
transition-producing computation can organize its actual forward execution
through stable learned operations under explicit state-conditioned control over
their contributions. We test that proposition constructively in autoregressive
language modeling, not the novelty of any constituent ingredient.

Architectural exposure gives operation identities, applicability coefficients,
operation responses, and controlled contributions direct measurement and
intervention targets. In the tested protocol, selected contributions are
causally accessible through suppression and exact restoration. This does not
establish semantic purity, universal cross-context meaning, arbitrary operation
replacement, control transplant, or program synthesis.

\DAWN{} is one constructive realization. It uses persistent rank-1 RW operators,
continuous-valued applicability in a \(d_c=256\) operation-address space, and
normalized additive composition across Q/K/V and the residual-state transition
(RST). Applicability is controlled through this low-dimensional representation,
whereas each applicable RW operation still reads the original
\(D\)-dimensional representation and writes in \(\R^D\). Thus low-dimensional
applicability control is not a low-dimensional content read. The evaluated
implementation considers the full type-specific operator pool and executes RW
contractions densely. Its spherical operation-address geometry is a control
inductive bias. Neither geometry nor rank-1 structure, persistence, global
candidate visibility, or additive composition is required by the formulation.

The empirical study tests this realization through three ordered claims:
viability, conditional organization, and causal accessibility. Fail-closed
results delimit where the mechanistic localization procedure does not succeed;
efficiency is not a headline claim.

Our contributions are:
\begin{enumerate}
    \item \textbf{An organizational formulation for explicit learned
    state-transition computation and a constructive architectural test.} We
    connect operation-oriented program execution, continuous state-dependent
    dynamics, and learned modular computation as complementary precedents; use
    \((\mathcal O,C,\Gamma)\) as minimal organizational coordinates; formulate
    the State-Conditioned Program as one operation-oriented organization of these
    roles in actual forward execution; and test the resulting architectural
    hypothesis through \DAWN{} in autoregressive language modeling. \DAWN{}
    provides one concrete realization using persistent rank-1 RW operations and
    low-dimensional continuous geometric applicability.
    \item \textbf{A constructive empirical test with bounded evidence.} We test
    viability at 400M scale, report descriptive results for the 1.3B \DAWN{}
    checkpoint (1.272B parameters), provide evidence of conditional organization
    through state-conditioned exact applicability support, and test causal
    accessibility through a discovery-frozen IOI circuit confirmed by
    suppression, matched controls, and exact restoration. Fail-closed results
    bound where the same localization procedure did not succeed.
\end{enumerate}

\section{Related Work}
\label{sec:related}

\subsection{Operational views of computation and state transition}

Structural operational semantics describes computation through transitions
between configurations governed by rules \cite{plotkin1981structural}. We use
only the associated operation-oriented view: execution can be organized around
identifiable operations and conditions under which they apply. \DAWN{} does not
define programming-language syntax, guarded commands, inference rules, loops, or
formal operational semantics. Its overlap with dynamics begins at the abstract
substrate \(x\mapsto x'\), not at a claimed formal equivalence.

\subsection{Continuous state-conditioned dynamics and local-model composition}

Dynamical modeling supplies the primary mathematical lineage for conditional
transition organization. Fuzzy sets formalized continuously graded membership
rather than Boolean set inclusion \cite{zadeh1965fuzzy}. Takagi--Sugeno models
couple graded premise validity to local consequent models and aggregate their
effects \cite{takagi1985fuzzy}. Local model networks likewise interpolate simple
component models, each with a limited operating regime, into a global nonlinear
model \cite{johansen1995identification}. Gain scheduling and linear
parameter-varying formulations organize families of local laws or controllers by
measured scheduling variables, while also making clear that properties of local
designs do not automatically transfer to the varying global system
\cite{rugh2000research}.

Learned-dynamics research also makes transition structure an explicit learned
object. Decomposed Linear Dynamical Systems learn dynamics components whose
sparse combinations represent overlapping nonstationary dynamics
\cite{mudrik2024decomposed}; Neural Hybrid Automata learn discrete modes and
events with multi-mode continuous dynamics \cite{poli2021neural}. Neither is a
State-Conditioned Program or equivalent to \DAWN{}; both provide precedent for
explicit learned transition structure.

Together these traditions establish precedent for continuous state- or
regime-dependent validity, overlapping component transition laws, and weighted
effect composition. They are not mathematically identical to one another or to
\DAWN{}, and they need not expose persistent learned operations or causal
intervention interfaces. State-conditioned dynamics, continuous membership, and
weighted combinations of local models are therefore not novelty claims here.

\subsection{Learned conditional and modular computation}

Mixture-of-experts (MoE) computation already has the generic form
\begin{equation}
F(x)=\sum_i G_i(x)E_i(x),
\label{eq:moe_precedent}
\end{equation}
with input-dependent control over learned transformations. Sparsely-gated MoE
\cite{shazeer2017outrageously}, Switch Transformer \cite{fedus2022switch},
ModuleFormer \cite{shen2023moduleformer}, and CondConv
\cite{yang2019condconv} establish conditional routing and weighted expert
combination. PEER retrieves million-scale tiny experts \cite{he2024mixture},
while OmniMoE uses vector-level Atomic Experts and structured routing
\cite{shi2026omnimoe}. Continuous or differentiable control over learned
transformations is therefore precedent; MoE-style conditional computation and
continuous applicability are not claimed as novel.

Shared parameter identity across depth is likewise established: Universal
Transformers recur over shared layers \cite{dehghani2019universal}, MoEUT shares
expertized feed-forward and attention computation \cite{csordas2024moeut}, MoUE
uses a layer-agnostic expert pool \cite{chen2026mixture}, and UniPool exposes a
global pool through per-layer routers \cite{huang2026unipool}. SwitchHead routes
attention-side experts \cite{csordas2024switchhead}, UMoE shares experts across
attention and FFN computation \cite{yang2025umoe}, and Compositional Attention
separates query--key search from value retrieval
\cite{mittal2022compositional}. These architectures make reuse possible and
measurable; stable or shared identity does not by itself establish broad
functional reuse.

Routing Networks combine a router with learned function blocks and recursively
choose functions from the current activation, producing input-specific dynamic
function composition \cite{rosenbaum2018routing}. Recursive Routing Networks
apply the same organizing principle to language understanding through a shared
set of learned functions and a meta-learner router
\cite{cases2019recursive}. These are direct precedents for
state/input-conditioned selection and composition of learned functions. Recurrent
Independent Mechanisms use groups with independent transition dynamics that
compete to update when relevant and communicate through attention
\cite{goyal2021rim}. The declarative/procedural factorization associated with the
SCOFF line separates object-state files from reusable schemata that prescribe
updates \cite{goyal2021factorizing}. Neural Production Systems learn
condition--action rules whose state- and entity-dependent matches update entity
state \cite{didolkar2021neural}, and Reusable Slotwise Mechanisms dynamically
select reusable mechanisms to predict future object-slot states
\cite{nguyen2023reusable}.

These works expose learned mechanism identity and conditional use: RIM, SCOFF,
Neural Production Systems, and Reusable Slotwise Mechanisms emphasize world or
object-state dynamics, while Routing Networks and related executors emphasize
neural-computation functions. This is emphasis, not a strict taxonomy or a claim
that they solve the same problem.

Neural Interpreters use learned types and signatures to control access to
reusable functions in a hyperspherical type space \cite{rahaman2021dynamic}, and
Neural Attentive Circuits separate a circuit generator from a modular executor
\cite{weiss2022neural}. These works establish conditional function access and
dynamic learned configurations. In this layer, both \(C_\theta\) and
\(O_{\theta,i}\) are learned objects: a local or component transition model can
be interpreted architecturally as a persistent learned computational operation.
Persistence supplies identity; whether that identity supports cross-context
reuse is an empirical question, not part of the definition.

\subsection{Read--write and low-rank operation structure}

Geva et al. interpret Transformer feed-forward layers as key-value memories with
input-sensitive coefficients over output directions
\cite{geva2021transformer}, and later decompose their additive updates into
parameter-vector sub-updates traced through vocabulary-space predictions
\cite{geva2022transformer}. Product Key Memory provides lineage for factorized
keys and structured retrieval \cite{lample2019large}, scalable memory layers
extend content-addressable trainable memories \cite{berges2024memory}, Fast
Weight Programmers provide an outer-product programming precedent
\cite{schlag2021linear}, and MoRAM uses rank-1 associative-memory experts with
intrinsic keys \cite{lu2026little}. Thus memory, retrieval, outer-product
structure, read/write decomposition, and rank-1 experts are precedents. They
provide the concrete transformation structure, not the general State-Conditioned
Program definition. \DAWN{} chooses
\(O_i(x)=\bar w_i(\bar r_i^\top x)\), then makes \(c_i(x)O_i(x)\) an explicit
controlled transition contribution under a separate applicability profile.

Implementation-level overlap in rank-1 experts, adaptive thresholds, geometric or
low-dimensional control, read/write structure, or causal intervention does not
imply shared conceptual ancestry, interpretation, or research program.

\subsection{Neural programming as adjacent precedent}

Neural Programmer-Interpreters execute programs from supervised traces
\cite{reed2016npi}, Neural Programmer differentiably selects operations and data
\cite{neelakantan2016neuralprogrammer}, and Fast Weight Programmers connect
outer-product computation to instructions \cite{schlag2021linear}. RASP gives a
programming-language-level model of Transformer computation and realizability
\cite{weiss2021thinking}. Neural programs, differentiable operation selection,
and program-level descriptions of Transformer computation therefore have
precedent. The State-Conditioned Program analogy is deliberately weaker and
different: it concerns the organization of conditional learned state-transition
operations, with no requirement for symbolic source, discrete instruction
traces, human-readable predicates, named primitive semantics, or formal
structural operational semantics.

\subsection{Causal intervention as evaluation methodology}

Mechanistic interpretability localizes components through causal intervention.
IOI work isolates an attention-head circuit \cite{wang2023interpretability},
AtP* approximates activation patching at scale \cite{kramar2024atp}, and RAVEL
evaluates localization against multiple causal criteria
\cite{huang2024ravel}. This literature distinguishes candidate localization from
causal confirmation. \DAWN{} supplies controlled RW-operation contributions as
intervention targets; its discovery freeze, suppression, matched controls, and
exact restoration test the specified IOI circuit without implying general
semantic interpretability. This work supplies a validation methodology, not a
constructive ingredient of the architecture.

The traditions above are neither historically linear nor mathematically
identical. Program and execution accounts contribute an operation-oriented view
of conditioned computation. Dynamics and control provide the mathematical
precedent for continuously state-dependent validity and composition of transition
effects. Learned modular systems show that control and computational components
can themselves be learned and assigned persistent identities. RW/fast-weight
work supplies elementary transformation structures, while causal-intervention
work provides methodology for testing whether exposed constituents matter to
behavior. We connect these precedents at the level of learned state-transition
computation itself: whether execution can be organized so that conditions,
operations, and transition effects become explicit computational objects.

At the coordinate level, one family of such organizations can be expressed as
\[
F_\theta(x)=
\Gamma\!\left(\{c_{\theta,i}(x)O_{\theta,i}(x)\}_i\right),
\]
using \((\mathcal O,C,\Gamma)\) as minimal organizational coordinates. This is an
organizational form, not the architectural hypothesis and not a universal theory
of learned dynamics. A State-Conditioned Program is one program-oriented
formulation in which these roles participate in actual forward execution. The
architectural hypothesis is the proposition that learned transition-producing
computation can realize that organization through stable learned operations under
explicit state-conditioned control over their contributions. \DAWN{} is one
concrete realization, at the granularity of a controlled RW contribution,
\[
c_i(x)O_i(x)=c_i(x)\,\bar w_i\big(\bar r_i^{\top}x\big).
\]
Applicability determines when and how much the operation contributes, the read
direction extracts a state-dependent scalar from the current representation,
and the write direction determines the resulting transition effect. The present
constructive test domain is autoregressive language modeling; the broader problem
is not limited to that domain.

\section{DAWN-SRW}
\label{sec:method}

\subsection{State Transitions and Control--Operation Factorization}
\label{sec:state-transitions}
\label{sec:transition-hypothesis}

Let
\[
X_\ell=(x_{\ell,1},\ldots,x_{\ell,T})\in\R^{T\times D}
\]
denote the sequence representation at depth \(\ell\). A Transformer-like block
produces state transitions through attention and a subsequent residual transformation:
\begin{align}
Y_\ell&=X_\ell+\operatorname{Attn}_\ell(X_\ell),\\
X_{\ell+1}&=Y_\ell+R_\ell(Y_\ell).
\label{eq:state_transitions}
\end{align}
Equations~\ref{eq:organizational_coordinates}--\ref{eq:general_composition}
provide a coordinate-level organizational form, not a sufficient condition for
computational organization or the architectural hypothesis itself. That
hypothesis is that learned transition-producing computation can organize its
actual forward execution through stable learned operations under explicit
state-conditioned control over their contributions. \DAWN{} exposes one
realization through persistent RW operators, an explicit applicability profile,
and normalized additive composition. For a token state \(x\in\R^D\) and
computation type \(a\in\{Q,K,V,\mathrm{RST}\}\), define
\begin{equation}
F_a:\R^D\rightarrow\R^D.
\label{eq:composed_output}
\end{equation}
Here \(F_Q,F_K,F_V\) parameterize relational interaction and
\(F_{\mathrm{RST}}\) parameterizes a direct residual-state update. Rank-1 RW
operators, operation-address geometry, persistent identities, full candidate
scope, dense physical execution, and normalized additive composition are
realization choices, not requirements of either the formulation or the
architectural hypothesis. We use \emph{operation} for the abstract
transformation, \emph{operator} for an indexed learned architectural object
implementing an operation, and \emph{RW operator} for the rank-1 realization
used here. Thus the computational formulation, architectural hypothesis, DAWN
realization, and geometry implementation are distinct levels of claim.

\begin{figure*}[t]
\centering
\setlength{\fboxsep}{4pt}
\small
\textbf{Language modeling as representation-state transitions}\\[0.35em]
\begin{tabular}{c@{\hspace{0.12em}\(\longrightarrow\)\hspace{0.12em}}c@{\hspace{0.12em}\(\longrightarrow\)\hspace{0.12em}}c@{\hspace{0.12em}\(\longrightarrow\)\hspace{0.12em}}c@{\hspace{0.12em}\(\longrightarrow\)\hspace{0.12em}}c}
\fbox{\parbox[c][0.92cm][c]{0.14\textwidth}{\centering
\textbf{Representation state}\\\(X_\ell\)}} &
\fbox{\parbox[c][0.92cm][c]{0.17\textwidth}{\centering
\textbf{Relational interaction}\\attention via \(F_Q,F_K,F_V\)}} &
\fbox{\parbox[c][0.92cm][c]{0.14\textwidth}{\centering
\textbf{Contextual representation}\\\(Y_\ell\)}} &
\fbox{\parbox[c][0.92cm][c]{0.16\textwidth}{\centering
\textbf{Residual-state transition}\\via \(F_{\mathrm{RST}}\)}} &
\fbox{\parbox[c][0.92cm][c]{0.10\textwidth}{\centering
\textbf{Next state}\\\(X_{\ell+1}\)}}
\end{tabular}\\[-0.05em]
\begin{tabular}{cc}
\parbox{0.34\textwidth}{\centering
\(F_Q,F_K,F_V\) parameterize relational interaction\\\(\Downarrow\)} &
\parbox{0.34\textwidth}{\centering
\(F_{\mathrm{RST}}\) directly updates the residual state\\\(\Downarrow\)}
\end{tabular}\\[0.15em]
\textbf{\DAWN{}: normalized additive composition}\\[0.35em]
\begin{tabular}{c@{\hspace{0.5em}}c@{\hspace{0.5em}}c@{\hspace{0.5em}}c@{\hspace{0.5em}}c}
&& \fbox{\parbox[c][0.68cm][c]{0.23\textwidth}{\centering
\textbf{Low-D applicability path}\\\scriptsize
\(x\!\to\!q_a(x)\!\to\!c_{a,i}(x)\)}} && \\[-0.28em]
\fbox{\parbox[c][0.68cm][c]{0.14\textwidth}{\centering
\textbf{Representation state}\\\(x\)}} & \(\nearrow\) && \(\searrow\) &
\fbox{\parbox[c][0.68cm][c]{0.18\textwidth}{\centering
\textbf{Composition result}\\\(F_a(x)\)}} \\[-0.28em]
& \(\searrow\) && \(\nearrow\) & \\[-0.28em]
&& \fbox{\parbox[c][0.86cm][c]{0.23\textwidth}{\centering
\textbf{Full-D RW response path}\\\scriptsize
\(x\!\to\!\bar r_{a,i}^{\top}x\!\to\!O_{a,i}(x)\)\\[-0.12em]
persistent pool; shared across depth}} &&
\end{tabular}\\[0.15em]
\caption{The two-level construction. Top: a language-model block produces state
transitions through relational interaction and a subsequent residual-state
transition. Bottom: the current state follows two distinct paths. The
applicability path uses the low-dimensional operation-address representation to
produce \(c_{a,i}(x)\); the RW response path reads a scalar
\(\bar r_{a,i}^{\top}x\) from the full \(D\)-dimensional representation and
writes \(O_{a,i}(x)\) along \(\bar w_{a,i}\). Their product is a controlled
operation contribution, and normalized additive composition produces
\(F_a(x)\). The type-specific operator pool is persistent across depth, candidate
scope is full, and physical RW execution is dense.}
\label{fig:architecture}
\end{figure*}

\subsection{Read-Write Operators}
\label{sec:rw}

Each RW operator owns learned read and write parameters. With
\(\operatorname{unit}(z)=z/\max(\norm{z}_2,\epsilon)\), define normalized
directions \(\bar r_{a,i}=\operatorname{unit}(r_{a,i})\) and
\(\bar w_{a,i}=\operatorname{unit}(w_{a,i})\). The RW operator executes
\begin{equation}
O_{a,i}(x)=\langle x,\bar r_{a,i}\rangle\bar w_{a,i}
=\bar w_{a,i}\bar r_{a,i}^{\top}x.
\label{eq:rw_operator}
\end{equation}
One controlled RW contribution can therefore be read as
\begin{equation}
\begin{aligned}
c_{a,i}(x)O_{a,i}(x)
={}&\underbrace{c_{a,i}(x)}_{\text{applicability}}
\underbrace{\big(\bar r_{a,i}^{\top}x\big)}_{\text{scalar state read}}\\[-0.1em]
&\times\underbrace{\bar w_{a,i}}_{\text{write direction}}.
\end{aligned}
\label{eq:rw_semantic_roles}
\end{equation}
Applicability and the read response are both state-dependent but serve different
roles. Applicability determines whether and to what degree the operation
contributes; the read direction extracts an operation-specific scalar operand
from the current representation; and the write direction determines how that
operand changes the state. Q/K/V RW operators compose attention interfaces; RST
RW operators compose a direct residual-state update.

\subsection{State-Conditioned Applicability}
\label{sec:applicability}

For each computation type, the current state induces an applicability profile
over its \(N_a\) persistent RW operators,
\begin{equation}
C_a(x)=\big(c_{a,1}(x),\ldots,c_{a,N_a}(x)\big),
\label{eq:applicability_profile}
\end{equation}
whose coordinates determine whether and to what degree each RW operator
participates. Section~\ref{sec:operation-space} defines the state-derived query
\(q_a(x)\) and operator address \(k_{a,i}\). Their applicability score and
applicability threshold are
\begin{align}
s_{a,i}(x)&=q_a(x)^\top k_{a,i},\\
\tau_a(x)&=-1+2\sigma\!\left(xu_a+\beta_a\right).
\label{eq:applicability_score_threshold}
\end{align}
The evaluated \texttt{linear\_angular} rule defines
\begin{equation}
c_{a,i}(x)=
\operatorname{clip}\!\left(
\frac{s_{a,i}(x)-\tau_a(x)}
{\max(1-\tau_a(x),10^{-4})},0,1\right).
\label{eq:applicability}
\end{equation}
We call \(c_{a,i}(x)\in[0,1]\) the \emph{continuous-valued applicability}
of operator \(i\) at state \(x\). Its exact applicability support is
\begin{equation}
\begin{aligned}
\mathcal{I}_a(x)
&=\{i:c_{a,i}(x)>0\}\\
&=\{i:s_{a,i}(x)>\tau_a(x)\}.
\end{aligned}
\label{eq:applicable_support}
\end{equation}
When \(s_{a,i}(x)\leq\tau_a(x)\), an operator contributes exactly zero; when
\(s_{a,i}(x)>\tau_a(x)\), \(0<c_{a,i}(x)\leq1\) varies continuously with the
applicability score. Thus the same applicability score and applicability
threshold specify exact non-participation and continuous-valued applicability.

\subsection{Normalized Additive Composition}
\label{sec:operator-composition}

\DAWN{} uses normalized additive composition as one concrete realization of the
composition rule \(\Gamma\). The rule combines controlled RW-operator
contributions; \(F_a(x)\) denotes the resulting output, not the rule itself.
With normalization factor \(Z_a(x)\),
\begin{align}
Z_a(x)&=\left[\max\!\left(\sum_i c_{a,i}(x),1\right)\right]^{p_a},\\
F_a(x)&=\frac{\sqrt{D/L}}{Z_a(x)}
\sum_i c_{a,i}(x)\bar w_{a,i}\bar r_{a,i}^{\top}x.
\label{eq:implemented_composition}
\end{align}
In each term, \(c_{a,i}(x)\) specifies state-conditioned applicability and
\(\bar w_{a,i}\bar r_{a,i}^{\top}x\) is the operation response; their product is
a controlled operation contribution. This additive form is a deliberate
realization choice: at a fixed state, it preserves an explicit operator-wise
decomposition of the transition-producing output, making individual controlled
contributions directly measurable and available as intervention targets; the
reported interventions suppress or restore selected numerator terms while
leaving the normalization factor unchanged. A nonlinear mixer would generally
introduce cross-operation interactions that obscure this exact decomposition;
because both \(c_{a,i}(x)\) and \(Z_a(x)\) are state-dependent, the overall mapping
\(x\mapsto F_a(x)\) remains generally nonlinear. Exact composition parameters
are listed in Appendix~\ref{app:equations}.

\subsection{Operation-Address Space}
\label{sec:operation-space}

In this realization, evaluating the condition under which an operation
contributes becomes a learned geometric relation between the current state and
that operator. The operation-address space is an experimentally controllable
inductive bias for state--operation applicability, not part of the
State-Conditioned Program definition. Representation and RW computation remain
in \(\R^D\), while the query and operator addresses
\(q_a(x),k_{a,i}\in\R^{d_c}\) are control variables. The control projection need
only preserve distinctions useful for applicability; once an operation is
applicable, Equation~\ref{eq:rw_operator} still reads the original
\(D\)-dimensional representation state. Two states can therefore have similar
applicability profiles yet different full-dimensional RW read responses and
transition contributions. This is an architectural consequence of the equations,
not an experimentally established equivalence relation. Here, \(d_c=256\) and
\begin{equation}
q_a(x)=\operatorname{unit}(xP_a+b_a)\in\R^{d_c}.
\label{eq:query}
\end{equation}
A shared probe pair \(P_r,P_w\in\R^{D\times d_c}\) maps functional read and
write directions into the operator address
\begin{equation}
k_{a,i}=\operatorname{unit}\!\left(
\operatorname{unit}(\bar r_{a,i}^{\top}P_r)\od
\operatorname{unit}(\bar w_{a,i}^{\top}P_w)\right)\in\R^{d_c}.
\label{eq:rw_address}
\end{equation}
Each operator address is a learned function of what its RW operator reads and
where it writes, rather than an independently parameterized address. The
operator address is not the scalar operand read by the RW operator: it is a
low-dimensional coordinate used to evaluate applicability of the
full-dimensional read--write transformation.

Because both \(q_a(x)\) and \(k_{a,i}\) are unit-normalized, they lie on the
unit sphere \(\mathbb{S}^{d_c-1}\). Writing \(\theta_{a,i}(x)\) for the angle
between them, \(s_{a,i}(x)=q_a(x)^\top k_{a,i}=\cos\!\big(\theta_{a,i}(x)\big)\),
so \(s_{a,i}(x)>\tau_a(x)\) is equivalent to
\(\theta_{a,i}(x)<\arccos\!\big(\tau_a(x)\big)\). Thus the exact applicability
support corresponds geometrically to a state-dependent spherical cap, or
angular neighborhood, centered at \(q_a(x)\) with angular radius
\(\arccos\!\big(\tau_a(x)\big)\). Within this neighborhood,
\(c_{a,i}(x)\) gives continuous-valued applicability that varies with the cosine
applicability score above the threshold and determines to what degree the
operator contributes. This is the spherical realization of
\((x,O_{a,i})\mapsto c_{a,i}(x)\); exact implementation details are restated in
Appendix~\ref{app:equations}. The construction induces locality in address
geometry, but the present measurements do not establish that geometric
proximity corresponds to functional or causal similarity.

\subsection{Attention-Mediated and Residual Transitions}
\label{sec:block}

For \(\widetilde X_\ell=\operatorname{LN}^{(1)}_\ell(X_\ell)\), three
normalized additive compositions produce the attention interfaces:
\begin{align}
Q_\ell&=F_Q(\widetilde X_\ell),\nonumber\\
K_\ell&=F_K(\widetilde X_\ell),\nonumber\\
V_\ell&=F_V(\widetilde X_\ell),\nonumber\\
A_\ell&=\operatorname{softmax}\!\left(
\frac{Q_\ell K_\ell^\top}{\sqrt{d_h}}+M_{\mathrm{causal}}\right)
\!,\nonumber\\
\Delta X_{\mathrm{attn},\ell}&=A_\ell V_\ell W^O_\ell,\\
Y_\ell&=X_\ell+\Delta X_{\mathrm{attn},\ell}.
\label{eq:attention}
\end{align}
Q and K use a shared type-specific operator pool, while V uses a separate
operator pool. A second
normalization \(\widetilde Y_\ell=\operatorname{LN}^{(2)}_\ell(Y_\ell)\) then
queries the RST operator pool:
\begin{equation}
X_{\ell+1}=Y_\ell+F_{\mathrm{RST}}(\widetilde Y_\ell).
\label{eq:rst}
\end{equation}
Thus RST applicability is evaluated from the contextualized post-attention
state, and its normalized additive composition produces the direct
residual-state update.

\subsection{Persistent Operator Identity and Shared Control}
\label{sec:reuse}

Type-specific operator pools and their state-conditioned control parameterizations
are shared across all \(L\) layers. Evolving normalized states therefore query
the same operator pools rather than a separately parameterized profile at each
depth. Persistence makes reuse measurable; it does not establish that broad
cross-context reuse emerges. The two
LayerNorms and attention output projection \(W^O_\ell\) remain layer-specific;
Appendix~\ref{app:settings} gives the complete sharing inventory.

\subsection{Candidate Scope and Physical Execution}
\label{sec:candidate-execution}

For computation type \(a\), the candidate scope is the full type-specific
operator pool. Thus
\[
\mathcal{I}_a(x)\subseteq\{1,\ldots,N_a\},
\]
with all \(N_a\) operators considered as candidates before applicability is
determined. Applicability support is generally a strict subset, while physical
RW execution is dense. Sharding, chunking, and accounting details appear in
Appendix~\ref{app:compute}.

\section{Experimental Setup}
\label{sec:setup}
\label{sec:protocol}

The primary model is \DAWN{} 400M (393.8M parameters, approximately 40B packed
C4 tokens), compared with a parameter- and training-token-matched dense
Transformer control (393.9M parameters). The 1.3B checkpoint, trained on approximately 20B
tokens, is descriptive only. Complete
architectural and training settings are in Appendix~\ref{app:settings}.

Performance evaluation uses packed C4 validation and six zero-shot tasks:
LAMBADA, HellaSwag, PIQA, ARC-Easy, ARC-Challenge, and WinoGrande. The two 400M
models use identical evaluation protocols; complete metrics, splits, and software
details are in Appendix~\ref{app:benchmarks}.

For causal analysis, IOI discovery ranks operator sites and freezes a circuit
before independent validation and held-out test. Suppression removes the frozen
sites' execution numerator, four matched-control families test specificity, and
exact restoration reinserts the selected same-example numerator under the
unchanged normalization factor. Appendix~\ref{app:ioi} gives the selection criteria, unit
counts, statistics, and complete intervention protocol.

\section{Results}
\label{sec:results}

\subsection{Language-model performance}

Table~\ref{tab:headline-performance} consolidates the headline results. At
400M scale, \DAWN{} differs from the matched dense control by
\(+0.004318918\) C4 cross-entropy and \(+0.003220495\) six-task mean. These
small measured gaps establish viability through performance retention, not
superiority or formal equivalence.

\begin{table}[t]
\centering
\caption{Headline language-model results. The 1.3B row is descriptive.}
\label{tab:headline-performance}
\small
\begin{tabular}{@{}lrrr@{}}
\toprule
Model & Parameters & C4 CE & Six-task mean\\
\midrule
\DAWN{} 400M & 393.8M & 2.940213 & 0.462872\\
Dense 400M & 393.9M & 2.935894 & 0.459652\\
\DAWN{} 1.3B & 1.272B & 2.835879 & 0.485844\\
\bottomrule
\end{tabular}
\end{table}

The descriptive 1.3B checkpoint reaches 2.835879 C4 cross-entropy and a
0.485844 six-task mean; without a matched dense control, this is not a scaling
comparison. Appendix~\ref{app:benchmarks} reports C4 accuracy and perplexity,
the complete task-level breakdown, and standard errors.

\subsection{State-Conditioned Exact Applicability Support}
\label{sec:support}

At all 9,994,240 packed-C4 positions, Q, K, V, and RST each have strict-subset
applicability support according to \(c_{a,i}(x)>0\) (Table~\ref{tab:support}).
This provides evidence of conditional organization through state-conditioned
exact applicability support under Equation~\ref{eq:applicability}, with
strict-subset rather than full-pool participation.

\begin{table}[t]
\centering
\caption{Exact applicability support according to \(c_{a,i}(x)>0\). Q/K union
counts an operator in the applicability support for Q or K at the same layer and
token.}
\label{tab:support}
\small
\begin{tabular}{lrrr}
\toprule
Computation & Pool & Mean support & Pool fraction\\
\midrule
Q & 4,600 & 487.896 & 10.6064\%\\
K & 4,600 & 386.722 & 8.4070\%\\
Q/K union & 4,600 & 812.388 & 17.6606\%\\
V & 14,500 & 1,649.994 & 11.3793\%\\
RST & 29,300 & 2,017.703 & 6.8864\%\\
\bottomrule
\end{tabular}
\end{table}

Restricting RW execution to exact applicability support would substantially
reduce static accounting, but the resulting estimate remains above the matched
dense control; this is not measured latency, energy, or hardware speedup. Detailed accounting and
Figure~\ref{fig:flops} are in Appendix~\ref{app:compute}.

\subsection{Frozen IOI circuit and causal intervention}
\label{sec:causal}

Discovery and confirmation are separated: a 4,096-site Q/K-centered circuit is
frozen before validation, and held-out test reuses the unchanged specification.
Complete discovery statistics are in Appendix~\ref{app:ioi}.

\begin{table*}[t]
\centering
\caption{Frozen IOI circuit results. Accuracy arrows show intact to
suppressed. Recovery is the fraction of the suppression effect recovered by
exact same-example numerator restoration.}
\label{tab:ioi}
\small
\resizebox{0.98\textwidth}{!}{
\begin{tabular}{lrrrrrr}
\toprule
Split & Units & Intact margin & Suppressed margin & Drop (95\% CI) &
Accuracy & Recovery / restored acc.\\
\midrule
Validation & 123 & 3.465456 & -4.399331 &
7.864787 [7.490080, 8.241263] &
1.000000 \(\rightarrow\) 0.008130 & 0.999661 / 0.991870\\
Held-out test & 114 & 3.254720 & -4.155613 &
7.410332 [7.016485, 7.818472] &
1.000000 \(\rightarrow\) 0.008772 & 0.998406 / 1.000000\\
\bottomrule
\end{tabular}
}
\end{table*}

Suppression causes large behavior changes on both splits and exceeds every
replicate in all four matched-control families. Activation-matched controls are
strongest; frozen-minus-control margin separation is 2.464297 on validation and
1.937754 on test (both \(p=0.00049975\)). Exact restoration recovers more than
99.8\% of the suppression effect. Together these interventions establish causal
accessibility for the evaluated IOI behavior; they do not imply a unique or
human-readable algorithm.

\subsection{Fail-Closed Results}
\label{sec:negatives}

Fail-closed results on RAVEL, ARC, and 1.3B IOI arise from distinct
preregistered criteria:
localization instability, failed causal suppression, and insufficient
cumulative importance, respectively. Exact statistics and decisions are in
Appendix~\ref{app:negative} and Table~\ref{tab:negative}.

\section{Discussion}
\label{sec:discussion}

\subsection{What the results establish}

The constructive test supports three ordered claims. First, the matched 400M
comparison establishes viability under the measured protocols. Second,
state-conditioned exact applicability support provides evidence of conditional
organization in the evaluated model; candidate scope is the full type-specific
operator pool and physical RW execution is dense. Third, the discovery-frozen
IOI circuit establishes causal accessibility: suppression changes behavior,
matched controls do not reproduce the effect, and exact same-example restoration
recovers it on independent validation and held-out data. Fail-closed results on
RAVEL, ARC, and 1.3B IOI delimit where this localization procedure did not
succeed.

These findings support one explicit realization under the evaluated conditions,
not an intrinsic or unique decomposition. Architectural exposure makes operator
identities, applicability coefficients, RW responses, and controlled
contributions direct measurement and intervention targets. The specified
suppression/restoration result establishes causal accessibility for the selected
IOI contributions; it does not establish human-readable operator semantics,
global modularity, a complete or unique causal decomposition, or transferable
program fragments.

Conceptually, the State-Conditioned Program is one program-oriented formulation
of the broader problem of explicit learned state-transition organization; its
continuous-composition mathematics has strong precedent in dynamical systems,
without implying formal operational semantics or mathematical equivalence.

\subsection{From the Computational Question to an Open Design Space}

Making learned state-transition organization explicit exposes four design axes:
control and applicability, operation and identity, coupling and composition, and
candidate scope and physical execution. The State-Conditioned Program supplies
the current program-oriented formulation, and \DAWN{} supplies one tested
point---rank-1 RW operators, an operation-address space, cross-depth persistence,
full candidate scope, dense execution, and normalized additive
composition---rather than requirements for the broader design space.

\paragraph{Control and applicability.}
How should state--operation applicability be represented? \DAWN{} derives
\(C_a(x)\) in a learned \(d_c=256\) operation-address space. Alternatives could
use functional parameters, intrinsic read activation, separate control
representations, or no geometry. The evaluated spherical realization induces
state-dependent angular neighborhoods, but it is unestablished whether this
geometry predicts functional or causal computation. A central open question is
whether learned control geometry recovers a geometry of transition requirements
rather than merely a convenient routing substrate.

\paragraph{Operation and identity.}
What should constitute a reusable learned operation? Rank-1 RW operators are
one granularity; alternatives could be finer, coarser, structured, or
hierarchical. \DAWN{} persists type-specific identities across depth, while other
scopes could be layer-local or recurrent. Persistence makes reuse measurable; it
does not establish that broad reuse emerges, guarantee a single function, or
imply semantic purity. Operations may remain distributed or polysemantic.

\paragraph{Coupling and composition.}
After control and operations are made explicit, how should their controlled
contributions combine into a transition? \DAWN{} derives addresses from
read/write parameters without identifying control variables with RW operations,
multiplies each response by its applicability, and uses normalized additive
composition partly to preserve exact operator-wise terms. Interactive or
nonlinear composition could trade separability for cross-operation structure;
sequential, hierarchical, recurrent, or more tightly coupled designs remain
possible. Whether an exactly additive decomposition is necessary for useful
causal access is itself an open question.

\paragraph{Candidate scope and physical execution.}
How much of an operator pool should candidate scope include? \DAWN{} uses the full
type-specific operator pool. Restricted or state-dependent retrieval, sparse
candidate selection, and sparse physical RW execution pose distinct
compute/quality questions; none is established by exact applicability support.

\subsection{Structural Questions Exposed by the Organization}

These questions concern explicit learned transition organization generally; the
present \DAWN{} geometry and language-model evidence instantiate but do not settle
them.

\paragraph{Dimension.}
The architecture distinguishes the representation dimension \(D\), control
dimension \(d_c\), the local span of an RW update, and the full operator pool.
The evaluated model is viable with \(d_c=256\ll D\), while RW responses still
read and write in full \(D\). This does not establish that \(d_c\) is minimal or
sufficient, that useful local execution is generally low-dimensional, or that
computation-selection complexity is intrinsically separate from state-content
complexity. A concrete question is how much of the ambient representation space
the applicable write directions and the full learned population must span.

\paragraph{Regime and geometry.}
Do applicability profiles define overlapping computational regimes, and does
proximity in control geometry predict similar transition requirements or causal
effects? The present study measures exact applicability support but not regime
equivalence, program similarity, or geometry--function correspondence.

\paragraph{Reuse and composition.}
Persistent identities could become state-specific specialists, locally reused
operations, cross-context reusable operations, or recurring higher-order
motifs. Which organization emerges is unmeasured. It is likewise open whether
applicability depends primarily on what an operation reads, on the effect it
writes, or on their joint read--write identity; the present architecture uses
the joint construction but does not compare alternatives.

\paragraph{Causal invariance and counterfactual intervention.}
The present causal protocol validates suppression and exact same-example
restoration of one discovery-frozen circuit. Whether similarity in
state-conditioned operation organization predicts similar intervention effects,
and whether control-profile transplant, operation transfer, or counterfactual
transition execution is meaningful, remain future questions rather than
demonstrated capabilities.

\subsection{From Applicability Support to Physical Execution}
\label{sec:future-execution}

For a given representation state, many operators have exactly zero
applicability, but this does not establish restricted candidate scope or sparse
physical execution. In \DAWN{},
\[
\mathcal{I}_a(x)\subseteq\{1,\ldots,N_a\},
\]
so the candidate scope is the full type-specific operator pool and physical RW
execution is dense. Applicability support specifies which operators contribute;
candidate scope specifies which operators control considers; physical execution
specifies how computation maps to hardware. \textbf{Sparse
applicability support does not by itself imply hardware-exploitable sparse
computation.} The open problem is to restrict candidate scope and physical
execution while preserving language-model quality and conditional organization.

\subsection{Limitations}

Global candidate evaluation and dense RW execution remain expensive,
and performance may depend partly on the capacity of the full operator pool.
Even static accounting with RW execution restricted to exact applicability
support remains 2.1168 times the matched dense control's backbone. No
production latency, energy, or hardware advantage is measured.

The familiar state--condition--operation motif is not claimed as novel, and an
arbitrary post-hoc functional decomposition does not establish computational
organization. The study does not establish that the State-Conditioned Program
organization or its control--operation factorization is intrinsic, unique,
necessary, privileged, or optimal. Nor are the spherical geometry,
\(d_c=256\) control dimension, or correspondence between geometric and
functional organization shown to be intrinsic. In particular, the low-
dimensional applicability path is not shown to be minimal, sufficient, a
statistical quotient of the representation, or evidence that local execution is
generally low-dimensional; the geometry may instead be a high-capacity control
substrate for broadly distributed full-dimensional operations. Persistence of
operator identity likewise makes reuse measurable but does not establish broad
cross-context reuse.

Causal evidence is confined to the specified IOI setting. Fail-closed results on
RAVEL, ARC, and 1.3B IOI delimit mechanistic generalization; the differently
budgeted 400M and 1.3B checkpoints do not define a controlled scaling curve; and
individual operators need not map one-to-one to human-readable concepts.
Suppression/restoration establishes causal accessibility only for the specified
terms and behavior; it does not establish arbitrary program editing, semantic
interpretability, operator replacement, or control transplant.

\newpage
\balance
\section{Conclusion}

Learned state-transition computation is usually organizationally implicit. The
coordinates \(\mathcal P=(\mathcal O,C,\Gamma)\) provide a minimal description of
one family of explicit organizations, not a universal or necessary decomposition.
A State-Conditioned Program is one program-oriented formulation requiring these
roles to participate in actual forward execution. The architectural hypothesis
is that learned transition-producing computation can realize such organization
through stable learned operations under explicit state-conditioned control over
their contributions. \DAWN{} tests one concrete rank-1 RW realization in the
present constructive domain of autoregressive language modeling:
applicability is computed in \(d_c=256\), while each operation reads a scalar
from and writes a direction into the full \(D\)-dimensional representation.

The evidence is bounded: measured 400M viability, descriptive 1.3B results,
evidence of conditional organization from state-conditioned exact applicability
support, and causal accessibility through a
discovery-frozen IOI circuit tested by suppression, matched controls, and exact
restoration. Fail-closed results mark where localization did not succeed.

The main interpretation is that one learned high-dimensional state-transition
system can expose referable operation contributions while retaining viability,
not that its decomposition or geometry is intrinsic, unique, semantically pure,
or computationally optimal. The broader question is how learned state-transition
computation can retain the flexibility of continuous neural representations while
organizing the conditions, operations, and transition effects that constitute its
execution as explicit objects of measurement and intervention.

\clearpage
\appendix

\section{Exact Operation-Address and Composition Equations}
\label{app:equations}

This appendix restates the production computation as a compact reference. Raw
read and write parameters are normalized in float32 before being cast for
dense contractions:
\[
\bar r_i=\frac{r_i}{\max(\norm{r_i}_2,\epsilon)},\qquad
\bar w_i=\frac{w_i}{\max(\norm{w_i}_2,\epsilon)}.
\]
One pair of model-level probes is shared across QK, V, and RST. In canonical
notation, the operator address is
\[
k_{a,i}=\operatorname{unit}\!\left(
\operatorname{unit}(\bar r_{a,i}^{\top}P_r)\od
\operatorname{unit}(\bar w_{a,i}^{\top}P_w)\right).
\]
The query and applicability threshold use separate direct projections. The
\texttt{DirectTau} projection produces the implementation variable
\texttt{raw\_tau} \(=xu_a+\beta_a\), whose transformed value is the canonical
applicability threshold \(\tau_a(x)\):
\[
\begin{aligned}
q_a(x)&=\operatorname{unit}(xP_a+b_a),\\
\tau_a(x)&=-1+2\sigma(xu_a+\beta_a),&
s_{a,i}(x)&=q_a(x)^\top k_{a,i}.
\end{aligned}
\]
For the \texttt{linear\_angular} applicability rule:
\[
c_{a,i}(x)=
\operatorname{clip}\!\left(
\frac{s_{a,i}(x)-\tau_a(x)}{\max(1-\tau_a(x),10^{-4})},0,1\right)
\]
and padding rows are masked to zero. This is the continuous-valued applicability
used in normalized additive composition. With normalization factor
\(Z_a(x)=\max(\sum_i c_{a,i}(x),1)^{p_a}\), the production composed output
before causal attention or residual addition is
\[
F_a(x)=
\sqrt{\frac{D}{L}}\,
\frac{\sum_i c_{a,i}(x)\langle x,\bar r_{a,i}\rangle\bar w_{a,i}}
{Z_a(x)},
\]
with \(p_Q=p_K=0.5\), \(p_V=1.0\), and \(p_{\mathrm{RST}}=1.2\).
Operator rows are sharded over the model axis; local numerators and local sums of
applicability values are reduced across shards to produce the global numerator
and normalization factor. Chunking changes the contraction schedule, not the
equation.

\paragraph{Local Jacobian and update-span consequence.}
Fix \(x\) and abbreviate \(\alpha=\sqrt{D/L}\), \(c_i=c_{a,i}(x)\),
\(u_i=\bar r_{a,i}^{\top}x\), \(Z=Z_a(x)\), and \(g_i=c_i u_i/Z\). On a
neighborhood with fixed exact applicability support, and at points away from
clipping, support, and normalization kinks, the realized normalized RW
composition has
\begin{align*}
\nabla_x g_i
&=Z^{-1}\!\left(c_i\bar r_{a,i}+u_i\nabla_x c_i\right)
-Z^{-2}c_i u_i\nabla_x Z,\\
J_{F_a}(x)
&=\alpha\!\sum_{i\in\mathcal I_a(x)}
\bar w_{a,i}\nabla_x g_i^{\top}.
\end{align*}
For
\(
\mathcal W_a(x)=
\operatorname{span}(\{\bar w_{a,i}\}_{i\in\mathcal I_a(x)})
\), consequently
\[
\operatorname{Im}J_{F_a}(x)\subseteq\mathcal W_a(x),\qquad
\operatorname{rank}J_{F_a}(x)\leq\dim\mathcal W_a(x).
\]
For a residual transition \(T(x)=x+F_a(x)\), the same containment applies to
\(J_T(x)-I\), not generally to \(J_T(x)\). This is a local consequence of the
evaluated normalized additive RW realization under the stated regularity
conditions, not a theorem about every State-Conditioned Program.

\section{Complete Model and Training Settings}
\label{app:settings}

\begin{table*}[!t]
\centering
\caption{Model and training specification. The 1.3B row is descriptive and
does not form a controlled scaling curve with the 400M rows.}
\label{tab:model-spec}
\small
\resizebox{0.99\textwidth}{!}{
\begin{tabular}{lrrrrrrrr}
\toprule
Model & Parameters & Tokens & \(D\) & Layers & Heads & QK / V / RST rows & Checkpoint\\
\midrule
\DAWN{} 400M & 393,800,708 & 39,999,504,384 & 2,304 & 18 & 36 &
4,600 / 14,500 / 29,300 & 76,293\\
Dense Transformer 400M & 393,896,960 & 39,999,504,384 & 1,280 & 18 & 20 &
-- & 76,293\\
\DAWN{} 1.3B & 1,272,038,404 & 19,999,981,568 & 4,096 & 24 & 32 &
8,534 / 26,904 / 54,352 & 87,193\\
\bottomrule
\end{tabular}
}
\end{table*}

\begin{table*}[!t]
\centering
\caption{Detailed settings from the registered training configurations and
restored checkpoint manifests.}
\scriptsize
\renewcommand{\arraystretch}{0.94}
\begin{tabular}{lrrr}
\toprule
Setting & \DAWN{} 400M & Dense 400M & \DAWN{} 1.3B\\
\midrule
Tokenizer vocabulary & 30,522 & 30,522 & 30,522\\
Residual width & 2,304 & 1,280 & 4,096\\
Layers & 18 & 18 & 24\\
Heads / head width & 36 / 64 & 20 / 64 & 32 / 128\\
Maximum sequence length & 512 & 512 & 512\\
Feed-forward width & -- & 5,120 & --\\
Control dimension \(d_c\) & 256 & -- & 256\\
QK / V / RST operator-pool size & 4,600 / 14,500 / 29,300 & -- &
8,534 / 26,904 / 54,352\\
QK / V / RST chunks & 1 / 1 / 4 & -- & 1 / 2 / 8\\
QK / V / RST normalizer exponent \(p_a\) & 0.5 / 1.0 / 1.2 & -- &
0.5 / 1.0 / 1.2\\
Initial applicability-support fraction targets & 0.05 / 0.03 / 0.01 & -- &
0.05 / 0.03 / 0.01\\
Training batch & 1,024 & 1,024 & 448\\
Mesh (data \(\times\) model) & \(16\times2\) & \(16\times2\) & \(16\times2\)\\
Learning rate & \(3\times10^{-4}\) & \(3\times10^{-4}\) & \(2\times10^{-4}\)\\
Warmup ratio & 0.08 & 0.08 & 0.06\\
Weight decay & 0.1 & 0.1 & 0.1\\
Applicability-threshold learning-rate multiplier & 0.001 & -- & 0.001\\
Dropout / \texttt{router\_dropout} & 0 / 0 & 0 / -- & 0 / 0\\
Gradient checkpointing & yes & yes & yes\\
Seed & 1 & 1 & 1\\
\bottomrule
\end{tabular}
\end{table*}

The 400M \DAWN{} checkpoint is the primary performance and mechanistic model.
The dense 400M checkpoint is only a matched performance control. The 1.3B
\DAWN{} checkpoint is used for descriptive performance and one IOI localization
replication that produced a fail-closed result. The still-incomplete dense 1.3B
run is excluded.

The QK, V, and RST operator pools; the two RW-to-address probes; the Q/K/V
query projection; the RST query projection; and their \texttt{DirectTau} projections
are shared across all layers. Q and K share read/write rows but retain separate
query slices, applicability thresholds, applicability values, and outputs. The two
LayerNorms and attention output projection \(W^O_\ell\) are layer-specific;
positional and token embeddings and the final normalization are model-level
parameters. Depth is not an explicit operator-pool index: it changes the representation
state presented to the shared control interface.

\section{Benchmark Metrics and Protocol Details}
\label{app:benchmarks}

\smallhead{C4.}
All reported C4 rows use the same packed validation artifact and
\texttt{packed\_c4\_validation\_v1} protocol. The capped view contains 19,531
sequences. Exactly 19,520 sequences form 610 complete global evaluation
batches, producing 9,974,720 valid next-token targets; 11 incomplete-tail
sequences are excluded before scoring.

\smallhead{Zero-shot.}
The protocol is \texttt{mamba\_table3\_zero\_shot\_v1}, using stock
\texttt{lm-eval==0.4.2}, global batch 32, maximum length 512, length buckets
64/128/256/512, and zero few-shot examples. Evaluated split sizes are 5,153
for LAMBADA, 10,042 for HellaSwag, 1,838 for PIQA, 2,376 for ARC-Easy, 1,172
for ARC-Challenge, and 1,267 for WinoGrande. Task, dataset, tokenizer, and
software settings are identical between the two 400M models.

\subsection{Detailed benchmark results}

\begin{table}[!t]
\centering
\caption{Packed C4 validation on the identical 9,974,720 target tokens.}
\label{tab:c4}
\small
\begin{tabular}{lrrr}
\toprule
Model & CE & PPL & Accuracy\\
\midrule
\DAWN{} 400M & 2.940213 & 18.919873 & 0.424912\\
Dense 400M & 2.935894 & 18.838336 & 0.427100\\
\DAWN{} 1.3B & 2.835879 & 17.045382 & 0.437099\\
\bottomrule
\end{tabular}
\end{table}

\begin{table*}[!t]
\centering
\caption{Zero-shot accuracy. Parentheses contain standard errors. HS, AE, and
AC denote HellaSwag, ARC-Easy, and ARC-Challenge.}
\label{tab:zero-shot}
\scriptsize
\resizebox{\textwidth}{!}{
\begin{tabular}{lrrrrrrr}
\toprule
Model & LAMBADA & HS & PIQA & AE & AC & WinoGrande & Mean\\
\midrule
\DAWN{} 400M &
0.352028 (0.006654) & 0.458375 (0.004972) & 0.712187 (0.010563) &
0.484007 (0.010255) & 0.256826 (0.012767) & 0.513812 (0.014047) & 0.462872\\
Dense 400M &
0.353192 (0.006659) & 0.457080 (0.004971) & 0.698041 (0.010712) &
0.475589 (0.010248) & 0.249147 (0.012639) & 0.524862 (0.014035) & 0.459652\\
\DAWN{} 1.3B &
0.395886 (0.006813) & 0.510157 (0.004989) & 0.723069 (0.010440) &
0.509680 (0.010258) & 0.268771 (0.012955) & 0.507498 (0.014051) & 0.485844\\
\bottomrule
\end{tabular}}
\end{table*}

LAMBADA perplexities are \(24.0774\pm0.8005\) for \DAWN{} 400M,
\(24.0042\pm0.8073\) for dense 400M, and \(17.6358\pm0.5390\) for
\DAWN{} 1.3B.

\section{IOI Localization and Intervention Details}
\label{app:ioi}

\subsection{Frozen discovery}

Discovery admits 119 paired-correct independent units and ranks 456,697 sites.
Split-half rank stability is 0.913978765. The frozen selection rule is the
smallest prefix satisfying cumulative absolute importance at least 0.70 and
split top-\(k\) overlap at least 0.95. The selected \(k=4096\) prefix has:
\begin{itemize}
    \item cumulative absolute importance 0.719248713;
    \item 3,993 shared split sites, overlap 0.974853516, and Jaccard
    0.950940700;
    \item Q/K/V/RST counts \(=1744/2350/2/0\);
    \item coverage of every layer from 0 through 17.
\end{itemize}
Validation and test do not recompute or alter this selection.

\begin{table*}[!t]
\centering
\caption{Control effects and frozen-minus-control separation. Each control
family contains 100 replicates. All paired permutation \(p\)-values are
0.0004997501 with 2,000 permutations.}
\label{tab:controls}
\scriptsize
\renewcommand{\arraystretch}{0.96}
\begin{tabular}{lrrrr}
\toprule
Control & Validation drop & Frozen minus control (95\% CI) &
Test drop & Frozen minus control (95\% CI)\\
\midrule
Equal-count & 0.070886 & 7.793902 [7.437851, 8.161572] &
0.080276 & 7.330057 [6.922408, 7.716879]\\
Same-layer & 0.039389 & 7.825399 [7.440629, 8.180288] &
0.041073 & 7.369260 [6.977447, 7.754539]\\
Activation-matched & 5.400490 & 2.464297 [2.076200, 2.832697] &
5.472579 & 1.937754 [1.564818, 2.319706]\\
Route-frequency & 0.219162 & 7.645626 [7.289512, 8.031625] &
0.242665 & 7.167668 [6.779836, 7.545537]\\
\bottomrule
\end{tabular}
\end{table*}

\subsection{Suppression, controls, and restoration}

Let \(T_{\mathrm{corr}}\) and \(T_{\mathrm{src}}\) denote the correct-answer and
source-answer token sets. The primary score is
\[
\mathcal{M}=\sum_{t\in T_{\mathrm{corr}}}\log p(t)
-\sum_{t\in T_{\mathrm{src}}}\log p(t).
\]
Suppression zeros the frozen sites' execution numerator at specified token
positions. It leaves applicability and the normalization factor unchanged.
Controls use the same 4,096-site count and exclude all circuit sites.
Activation matching conditions on layer, computation type, and the discovery mean
pre-scale operator-output-norm quantile. Restoration executes a distinct
intervention: after suppression, it adds back the exact selected same-example
numerator under the unchanged normalization factor.

On validation, restoration returns the margin to 3.462791, leaving residual
drop 0.002665. On held-out test, it returns the margin to 3.242905, leaving
residual drop 0.011815. The respective recovery fractions are 0.999661155
(95\% CI [0.998654544, 1.000584125]) and 0.998405645
(95\% CI [0.997323737, 0.999385266]).

\section{Fail-Closed Results}
\label{app:negative}

\begin{table*}[!t]
\centering
\caption{Fail-closed results. Later phases were not opened after a
failed preregistered criterion.}
\label{tab:negative}
\small
\begin{tabularx}{\textwidth}{l l Y Y}
\toprule
Study & Passed & Failed criterion & Decision\\
\midrule
RAVEL 400M & Behavioral eligibility & Variable stability:
0.628762 / 0.676243 / 0.660622 & No frozen circuit; no validation, transfer, or test\\
ARC 400M & Stable discovery; 4,096 Q/K sites; restoration &
Suppression effect wrong sign; control separation failed &
Reject causal replication; held-out test closed\\
IOI 1.3B & Behavioral eligibility; rank stability 0.907188 &
4,096-site cumulative importance 0.635144 \(<0.70\) &
No frozen circuit; validation and test closed\\
\bottomrule
\end{tabularx}
\end{table*}

\smallhead{RAVEL.}
Discovery uses 142 independent units. Rank stability is 0.628762 for
Continent, 0.676243 for Country, and 0.660622 for Language, below the
preregistered requirement. No prefix passes the complete discovery rule, no
variable circuit is frozen, and validation, transfer, and test are not opened.

\smallhead{ARC.}
Discovery stability is 0.89386 and a 4,096-site Q/K-only circuit is frozen.
On 51 validation units, intact margin is 1.735534 and suppressed margin is
2.819139, giving a frozen margin drop of \(-1.083605\) (95\% CI
[\(-1.472180\), \(-0.690348\)]). The effect has the wrong sign and fails
control separation. Restoration recovery is 1.004083 (95\% CI
[0.994587, 1.013652]), but restoration alone does not pass the conjunction.
Held-out test remains closed.

\smallhead{1.3B IOI.}
The scale replication admits 125 independent units and passes rank stability
at 0.907187749. The largest allowed 4,096-site prefix reaches only 0.635144
cumulative absolute importance, below the 0.70 criterion. No circuit is frozen;
validation and test remain closed.

\section{Compute Accounting}
\label{app:compute}

The implementation shards operator rows across the model mesh and processes
them in dense chunks. The 400M configuration uses one QK chunk, one V chunk,
and four RST chunks. Every query scores all local operator addresses, and every
operator row participates in dense read and write contractions before
cross-shard reduction. RW operators execute in the full residual dimension
\(D\); the \(d_c\)-dimensional operation-address space is used only for applicability.

The applicability-support analysis evaluates 9,994,240 packed-C4 positions and
observes applicability scores and applicability thresholds before RW
application. It stores aggregate statistics,
not per-token operator identities or vectors. Static accounting uses batch 1,
sequence length 512, and \(1\) MAC \(=2\) FLOPs. It counts major
matmul/einsum operations and omits softmax, normalization, applicability elementwise
arithmetic, embedding lookup, residual additions, and communication
collectives.

Static forward accounting is 4.741260116 TFLOPs for current dense \DAWN{},
0.818230250 TFLOPs for dense addressing plus RW application restricted to exact
applicability support, and 0.386547057 TFLOPs for the matched dense control.
The estimate restricted to exact applicability support is 17.2577\% of current
\DAWN{} accounting but
still 2.1168 times the dense Transformer. These are static arithmetic estimates,
not measured latency, energy, or hardware speedup.

\begin{figure}[!t]
\centering
\includegraphics[width=\linewidth]{figures/flops_stacked.pdf}
\caption{Static backbone accounting at batch 1 and sequence length 512.
Address scoring remains dense in the estimate restricted to exact applicability
support. Current dense RW execution also applies all full-dimensional RW rows. No
production latency claim is made.}
\label{fig:flops}
\end{figure}

\begin{table*}[!t]
\centering
\caption{Backbone static accounting. ``Exact applicability support'' retains full dense
address scoring and replaces only full-pool RW application with the measured
applicability support according to \(c_{a,i}(x)>0\).}
\scriptsize
\renewcommand{\arraystretch}{0.96}
\begin{tabular}{lrrr}
\toprule
Component & Current dense \DAWN{} & Dense address + exact RW & Dense Transformer\\
\midrule
Address scoring + query/threshold projections & 0.293742 & 0.293742 & --\\
RW application & 4.306187 & 0.383157 & --\\
Causal attention & 0.043487 & 0.043487 & 0.024159\\
Dense output / QKV / FFN maps & 0.097845 & 0.097845 & 0.362388\\
\midrule
Total TFLOPs & 4.741260 & 0.818230 & 0.386547\\
\bottomrule
\end{tabular}
\end{table*}

The estimate restricted to exact applicability support is mathematically
output-equivalent if the executor
can skip all and only zero-weight RW terms. It does not include indexed
address retrieval and does not demonstrate kernel parity. Instrumented
profiling time is excluded from latency claims because histogram collection is
part of the measured path.

\section{Reproducibility Identities}
\label{app:repro}

The experiments reported here use \DAWN{} implementation revision v4.1.7.2,
recorded under the repository model identity \texttt{spatial-r1-v4.1.7.2}; literal
repository filenames and artifact paths retain the historical shorthand
\texttt{v4172}.

\begin{table*}[!t]
\centering
\caption{Primary immutable identities. Branch names are omitted because they
move; exact checkpoint steps, hashes, and artifact locations are the binding
identities.}
\scriptsize
\begin{tabularx}{\textwidth}{>{\raggedright\arraybackslash}p{0.13\textwidth}
>{\raggedright\arraybackslash}p{0.43\textwidth} Y}
\toprule
Evidence & Exact identity & Repository source\\
\midrule
\DAWN{} 400M & model \texttt{spatial-r1-v4.1.7.2}; step 76,293;
checkpoint identity \texttt{a7ce8afc...f91c17}; config hash
\texttt{08733ae4...5039b} &
\path{configs/train_config_v4172_400M_c4_40B_v4_64_ver1_den_qk0p5_v1p0_rst1p2.yaml}\\
Dense 400M & model \texttt{baseline-JAX}; step 76,293; config hash
\texttt{c3ce5d72...82fc25} &
\path{configs/train_config_baseline_tpu_400M_c4_40B_v4_64_tp2.yaml}\\
\DAWN{} 1.3B & model \texttt{spatial-r1-v4.1.7.2}; step 87,193;
checkpoint identity \texttt{2e38b41f...bf5f9}; config hash
\texttt{1e877263...4a5222} &
\path{configs/train_config_v4172_1p3B_c4_20B_v4_64_ver1_den_qk0p5_v1p0_rst1p2.yaml}\\
IOI circuit & parsed frozen-spec hash \texttt{6ec26618...418c};
site hash \texttt{573cfff8...0200}; selected-row hash
\texttt{1439f6f3...fe0c1} &
\path{configs/paper_ioi_frozen_circuit_v4172_400m.yaml}\\
Compute support & support hash \texttt{b1eac5f5...bdc2};
accounting hash \texttt{d7ad6115...984e} &
\path{assets/paper/v4172_400m_compute_support_step76293/}\\
\bottomrule
\end{tabularx}
\end{table*}

The machine-readable checkpoint registry is
\texttt{configs/paper\_checkpoint\_registry.yaml}. Exact evaluation commits,
software versions, data hashes, result roots, and complete artifact hashes are
recorded in the repository experiment ledger. Historical experiment branches
are evidence identities only; this source is maintained on the canonical main
branch.

\section*{Acknowledgments}

The author thanks Anthropic Claude, Claude Code, OpenAI ChatGPT, and OpenAI
Codex for assistance throughout this research. This work was developed with
extensive use of AI systems for literature exploration, conceptual exploration,
cross-domain framing, abstraction and terminology refinement, implementation
support, debugging, and manuscript preparation. In particular, these systems
were used as interactive assistants for refining the learned state-transition
formulation and the exposition of \DAWN{}. The research direction,
architectural decisions, experimental design, experiments, analyses,
interpretations, and final claims remain the responsibility of the author.

This research was supported by the Google TPU Research Cloud (TRC) program,
which provided computational resources for training and analysis.

\begingroup
\footnotesize
\setlength{\bibsep}{0pt}
\bibliography{references}
\bibliographystyle{icml2026}
\endgroup

\end{document}
